%\documentclass[
%sigconf,
% authordraft
%]{acmart}
\documentclass[sigconf, screen, review, anonymous]{acmart}
% NOTE that a single column version may required for submission and peer review. This can be done by changing the \doucmentclass[...]{acmart} in this template to
% \documentclass[manuscript,screen]{acmart}

% To ensure 100% compatibility, please check the white list of approved LaTeX packages to be used with the Master Article Template at
% https://www.acm.org/publications/taps/whitelist-of-latex-packages 
% before creating your document. The white list page provides information on how to submit additional LaTeX packages for review and adoption. Fonts used in the template cannot be substituted; margin adjustments are not allowed.
\RequirePackage[
    datamodel=acmdatamodel,
    style=acmnumeric,
    sorting=none
]{biblatex}
\newcommand{\ie}{\emph{i.e.}, }
\newcommand{\eg}{\emph{e.g.}, }
\newcommand\blfootnote[1]{%
  \begingroup
  \renewcommand\thefootnote{}\footnote{#1}%
  \addtocounter{footnote}{-1}%
  \endgroup
}

% \addbibresource{references.bib}

% \BibTeX command to typeset BibTeX logo in the docs
\AtBeginDocument{%
  \providecommand\BibTeX{{%
    \normalfont B\kern-0.5em{\scshape i\kern-0.25em b}\kern-0.8em\TeX}}}

% complete the rights form.
% \setcopyright{acmlicensed}
% \copyrightyear{2024}
% \acmYear{2024}
% \acmDOI{XXXXXXX.XXXXXXX}
\settopmatter{printacmref=false} % Removes citation information below abstract
% \renewcommand\footnotetextcopyrightpermission[1]{} % removes footnote with conference information in first column
% \pagestyle{plain} % removes running headers 
% \setcopyright{none}

% These commands are for a PROCEEDINGS abstract or paper.
% \acmConference[Conference acronym 'XX]{June 03--05, 2018}{Woodstock, NY}
%\acmBooktitle{Woodstock '18: ACM Symposium on Neural Gaze Detection, June 03--05, 2018, Woodstock, NY} 
% \acmISBN{978-1-4503-XXXX-X/18/06}


% Submission ID.
% \acmSubmissionID{xxxx}

% \citestyle{authoryear}

\begin{document}

%%
%% The "title" command has an optional parameter,
%% allowing the author to define a "short title" to be used in page headers.
\title[PointStream]{PointStream: A Generative AI and Pose Estimation Framework for Video Streaming}

\renewcommand{\shortauthors}{Artioli et al.}

\author{Emanuele Artioli}
\email{emanuele.artioli@aau.at}
\orcid{1234-5678-9012}
\affiliation{
  \institution{Alpen-Adria-Universitaet}
  \streetaddress{Universitaetstrasse 3}
  \city{Klagenfurt}
  \state{Kaernten}
  \country{Austria}
  \postcode{9020}
}
\author{Daniele Lorenzi}
\email{daniele.lorenzi@aau.at}
% \orcid{1234-5678-9012}
\affiliation{
  \institution{Alpen-Adria-Universitaet}
  \streetaddress{Universitaetstrasse 3}
  \city{Klagenfurt}
  \state{Kaernten}
  \country{Austria}
  \postcode{9020}
}
\author{Farzad Tashtarian}
\email{farzad.tashtarian@aau.at}
\orcid{1234-5678-9012}
\affiliation{
  \institution{Alpen-Adria-Universitaet}
  \streetaddress{Universitaetstrasse 3}
  \city{Klagenfurt}
  \state{Kaernten}
  \country{Austria}
  \postcode{9020}
}
\author{Christian Timmerer}
\email{christian.timmerer@aau.at}
\orcid{1234-5678-9012}
\affiliation{
  \institution{Alpen-Adria-Universitaet}
  \streetaddress{Universitaetstrasse 3}
  \city{Klagenfurt}
  \state{Kaernten}
  \country{Austria}
  \postcode{9020}
}

\begin{abstract}
Streaming high-quality video requires significant bandwidth, making it challenging to deliver smooth, low-latency experiences under network constraints. We propose PointStream, a novel approach that minimizes transmitted data by extracting and encoding dynamic scene elements while reconstructing the full visual content at the client side. Our method segments moving objects from a static background, tracks their motion, and transmits a compact representation consisting of object positions, poses, and minimal visual data. A generative model at the client side reconstructs high-fidelity frames using this encoded information, significantly reducing bandwidth consumption while preserving visual continuity. We demonstrate PointStream in the context of sports video streaming, where players are automatically segmented, tracked, and reconstructed. Our evaluation shows that PointStream maintains perceptual quality while achieving substantial bandwidth savings compared to traditional video compression techniques.
\end{abstract}

% The code below is generated by the tool at http://dl.acm.org/ccs.cfm.
\begin{CCSXML}
<ccs2012>
   <concept>
       <concept_id>10010147.10010178.10010224.10010245.10010252</concept_id>
       <concept_desc>Computing methodologies~Object identification</concept_desc>
       <concept_significance>500</concept_significance>
       </concept>
   <concept>
       <concept_id>10010147.10010178</concept_id>
       <concept_desc>Computing methodologies~Artificial intelligence</concept_desc>
       <concept_significance>500</concept_significance>
       </concept>
   <concept>
       <concept_id>10002951.10003227.10003241.10010843</concept_id>
       <concept_desc>Information systems~Online analytical processing</concept_desc>
       <concept_significance>300</concept_significance>
       </concept>
   <concept>
       <concept_id>10002951.10003227.10003251.10003255</concept_id>
       <concept_desc>Information systems~Multimedia streaming</concept_desc>
       <concept_significance>500</concept_significance>
       </concept>
   <concept>
       <concept_id>10010147.10011777.10011778</concept_id>
       <concept_desc>Computing methodologies~Concurrent algorithms</concept_desc>
       <concept_significance>300</concept_significance>
       </concept>
   <concept>
       <concept_id>10010147.10010371.10010395</concept_id>
       <concept_desc>Computing methodologies~Image compression</concept_desc>
       <concept_significance>500</concept_significance>
       </concept>
 </ccs2012>
\end{CCSXML}

\ccsdesc[500]{Computing methodologies~Object identification}
\ccsdesc[500]{Computing methodologies~Artificial intelligence}
\ccsdesc[300]{Information systems~Online analytical processing}
\ccsdesc[500]{Information systems~Multimedia streaming}
\ccsdesc[300]{Computing methodologies~Concurrent algorithms}
\ccsdesc[500]{Computing methodologies~Image compression}

\keywords{HTTP adaptive streaming, Generative AI, End-to-end architecture, Quality of Experience}

% \received{20 February 2007}
% \received[revised]{12 March 2009}
% \received[accepted]{5 June 2009}

\maketitle

\newpage
\section{Introduction}
\label{sec:introduction}

% \begin{figure}[t]
%     \centering
%     \includegraphics[width=\columnwidth]{figures/PointStreamOverview.pdf}    
%     \caption{Overview of the PointStream pipeline.}
%     \label{fig:overview}
% \end{figure}

\begin{figure*}
    \centering
    \includegraphics[width=1\linewidth]{figures/PS-overview.pdf}
    \caption{Enter Caption}
    \label{fig:enter-label}
\end{figure*}

In the digital age, video streaming has become ubiquitous, underpinning entertainment, education, and communication platforms. Traditional video streaming methods involve transmitting entire frames at regular intervals, encoding every pixel regardless of its relevance. While compression techniques like motion estimation in codecs such as H.264~\cite{h264_overview} and AV1~\cite{av1_overview} reduce redundancy, they still rely on full-frame encoding, leading to inefficiencies—particularly the repeated transmission of static background information. As video resolutions escalate to 4K and beyond, these inefficiencies become more pronounced, demanding substantial bandwidth and storage resources.

In contrast, the human visual system processes scenes with remarkable efficiency. Rather than analyzing every detail, the brain focuses on key elements such as moving objects and salient features, while peripheral vision provides contextual information. This selective attention allows for rapid scene recognition and understanding. Studies have shown that peripheral vision contributes significantly to scene recognition, with the brain integrating low-resolution peripheral information to guide eye movements and focus on relevant details.  Furthermore, prior visual experience facilitates the rapid understanding of scenes, enabling the brain to reconstruct missing information based on expectations\cite{Central and peripheral vision for scene recognition: A neurocomputational modeling exploration, Making Sense of Real-World Scenes, What you see is what you expect: rapid scene understanding benefits from prior experience}.

Inspired by the efficiency of human vision, we propose PointStream, a novel video streaming approach that prioritizes the transmission of essential dynamic information. Instead of sending full frames, PointStream identifies and isolates moving objects within a scene, encodes their motion and appearance using minimal data, and reconstructs them on the client side using generative models. Static background elements are transmitted once and reused, significantly reducing the data required for high-quality video playback. This method aligns with the human visual system’s strategy of focusing on dynamic elements and using contextual understanding to fill in details.

The primary contributions of this work include:
\begin{itemize}
    \item A system that emulates human visual efficiency by transmitting only essential motion data instead of full-frame images.
    \item A flexible architecture capable of integrating different detection, segmentation, pose estimation, and background stitching algorithms, ensuring adaptability to advancements in the field.
    \item A generative model to reconstruct high-fidelity video frames from minimal transmitted data, preserving perceptual quality while reducing bandwidth requirements.
\end{itemize}

\newpage

\section{Related Work}

When videos started being shared over the Internet, they were initially delivered via progressive download, i.e., downloaded in their entirety, and only then played, regardless of whether the viewer was interested in watching the video in its entirety~\cite{progressive}.
To overcome the shortcomings of progressive download, Adaptive Bitrate Streaming (ABR) was introduced. ABR detects a user’s bandwidth and buffer capacity in real-time, adjusting the quality of the media stream accordingly~\cite{ABR from Daniele? A survey paper?}. This approach minimizes buffering and provides a smoother viewing experience across varying network conditions.
Building upon ABR, Dynamic Adaptive Streaming over HTTP (DASH), also known as MPEG-DASH, and Apple HTTP Live Streaming (HLS), emerged as international standards~\cite{dash, hls}. Both segment multimedia files into smaller parts and delivers them over HTTP, allowing clients to adaptively choose segments that match their network conditions and device capabilities. This standardization facilitated broader compatibility across different devices and platforms.
As video resolutions and streaming demands continued to rise, more advanced codecs were developed to improve compression efficiency. Versatile Video Coding (VVC), also known as H.266, and AOMedia Video 1 (AV1) are currently regarded as state-of-the-art solutions, achieving up to 50\% bitrate savings over their predecessors at comparable quality~\cite{vvc, av1}. Despite their efficiency, both codecs require significantly higher computational complexity for encoding and decoding, making real-time applications more challenging and prompting exploration into different compression techniques. For example, artificial intelligence (AI) techniques have been integrated into video compression pipelines~\cite{ML for encoding, from Mohammad?}. AI-driven approaches optimize encoding parameters, predict optimal bitrate settings, and enhance overall compression performance, leading to better quality at reduced bitrates.
Leveraging AI even further, neural network-based compression methods have been developed. These approaches utilize deep learning models to learn compact representations of video data, achieving compression rates comparable to traditional methods, but still plagued by the same high computational requirement~\cite{some NeRV version from YiYing?}.
Another area of improvement for adaptive streaming consist of leveraging the computation power on the client side to improve video quality, with techniques such as frame interpolation and super-resolution~\cite{super resolution, frame interpolation, SRFI}. Such techniques, independent from network conditions, are a promising direction of research, and already well optimized for the quickly rising capabilities of modern GPU chips. However, they apply low-level (i.e., texture, edges, etc.) enhancements uniformly, without understanding the semantic content of the video.
Generative AI is poised to raise the ceiling in video streaming by understanding the content of a video, and performing targeted encoding and reconstruction based on high-level video features such as entire objects and complex movements. For example, ELVIS removes regions of a video, and inpaints them on the client side, primarily reducing the amount of background elements needed to be trasmitted over the network~\cite{elvis}. While innovative, this solution remains anchored to traditional block-based analysis and does not fully exploit object-based encoding.
This progression illustrates how each technological advancement addresses specific challenges in video streaming, at the same time introducing new considerations that necessitate further innovation. The landscape is therefore primed for PointStream, which aims to overcome these limitations by introducing a fully object-based encoding approach, enabling efficient and high-quality video streaming through intelligent scene reconstruction, in a similar fashion to the way the human brain parses visual information.

% Live video streaming has traditionally relied on two main pillars. 
% First, HTTP Adaptive Streaming (HAS) protocols use pre‐encoded representations~\cite{AVC, HEVC, AV1, VVC} to cope with network instability by dividing videos into segments and using adaptive bitrate (ABR) algorithms for quality selection~\cite{Wiegand et al. 2003, Sullivan et al. 2012}. 
% Second, content delivery networks (CDNs) and modern adaptive systems~\cite{Pensieve} have optimized client–server interactions to maintain user Quality of Experience (QoE). 
% However, despite these advances, bandwidth shortages due to network congestion and the high data demands of high-definition content persist.

% \subsection{Video Compression and Redundancy Removal}
% Conventional video codecs exploit spatial and temporal redundancies via techniques such as motion compensation and transform coding to achieve compression. 
% While these methods have matured over decades, they often require transmitting a substantial amount of data—even for scenes with largely static backgrounds. 
% % we could talk more about how blocks of a static background with objects moving on them, need to be re-encoded every time an object occludes them, since the codec does not know what is background and what is object.
% Recently, neural video codecs~\cite{NeRV, HNeRV} have emerged as alternatives that learn compact representations of video content. 
% Although promising, such approaches tend to shift the computational burden to the client and are not yet optimized for the real‐time constraints of live streaming.

% \subsection{Neural Video Restoration and Fidelity Reconstruction}
% Parallel to advancements in compression, deep learning–based video restoration techniques have been explored to enhance perceptual quality from low-bitrate or compressed streams. 
% Methods for frame interpolation, super-resolution, and inpainting~\cite{ELVIS} use convolutional neural networks and generative adversarial networks to recover high-quality visual details. 
% Yet, these approaches generally target the entire frame without explicitly exploiting the fact that many live streams exhibit predominantly static backgrounds. 
% In contrast, our method specifically targets scenarios where background elements remain largely unchanged while objects or people move—allowing the client to reconstruct high-fidelity static content using generative models, thereby reducing client load and transmitted data volume.

% \subsection{Content-Aware Adaptation in Live Streaming}
% A growing body of work has recognized that not all video segments are equally sensitive to quality degradation. 
% For instance, SENSEI~\cite{Zhang et al., NSDI '21} demonstrates that user sensitivity to quality changes varies dynamically with content—for example, rebuffering during a crucial moment in a live event is more detrimental than during routine play. 
% % this is why we base our segments into scenes instead of seconds.
% Similarly, recent systems like GRACE~\cite{Cheng et al., NSDI '24} leverage neural codecs to improve resilience against packet loss by reconstructing missing content on the fly. 
% These studies underscore the potential of incorporating semantic or content-aware insights into streaming systems. 
% Our work builds on this idea by concentrating on fidelity reconstruction of static scenes; by transmitting only essential motion information, we can devote computational resources at the client side (on modern GPUs) to reconstruct the crucial elements of a video with high visual fidelity.

\newpage

\section{Problem Description and Motivation}
To evaluate whether our approach brings meaningful improvements in video quality and computational efficiency, we first establish a baseline by measuring the performance of current state-of-the-art codecs under constrained bandwidth conditions.
The codec of choice is H.265, for its combination of adoption and encoding effiency.
As our reference codec, we select H.265 (HEVC) due to its widespread adoption and high compression efficiency. We encode the target video at 900 kbps and 960×540 resolution, following the HLS encoding ladder for adaptive streaming. This bitrate represents a moderate but constrained bandwidth scenario, reflecting real-world conditions where video quality often suffers due to network limitations. As shown in Table~\ref{table:motivation}, encoding the target video under these conditions results in a quality of X VMAF points, which may be insufficient for an optimal viewing experience.
Given the structured nature of tennis video segments, we explore a content-aware alternative that exploits the inherent redundancy in the scene. A typical broadcast consists of (i) A static background (court, audience, advertising boards), and (ii) a few dynamic elements (players, rackets, ball).
Instead of encoding all pixels in every frame, we propose a simple yet effective optimization:
\begin{itemize}
    \item Extract and store only the moving objects (players, rackets, and ball) as separate video streams.
    \item Compress and store the background as a single static image.
\end{itemize}
This approach ignores irrelevant motion (e.g., minor spectator movements, referees) while significantly reducing bitrate consumption—by up to X\%, as demonstrated in Table~\ref{table:motivation}, in a similar fashion to~\cite{elvis}. Despite these gains, this method still requires encoding full pixel-based representations of moving objects, leaving room for further improvements.
Encouraged by these findings, we take the next step toward a more efficient, object-aware approach to video streaming. Instead of merely isolating objects, we explore the potential of client-side computation to reconstruct complex motion and visual details from minimal transmitted data. This insight motivates the development of PointStream, a system designed to further reduce bitrate requirements while preserving and potentially even enhancing perceived video quality.

\newpage
\section{System Design}
PointStream is a system designed to efficiently transmit and reconstruct dynamic objects within scenes while minimizing the amount of data sent to the client. The system achieves this by extracting key object representations at the server and reconstructing them at the client using learned generative models. The design is structured around three core components: object extraction, background stitching, and client-side synthesis.

\subsection{PointStream - Server side}
The server-side processing is responsible for extracting dynamic objects from input frames. This is achieved through a combination of image segmentation, object tracking, and pose estimation techniques. These methods allow the system to isolate relevant objects, determine their spatial configurations, and track their evolution over time. The extracted representations contain minimal  data to enable accurate reconstruction on the client side. Additionally, the server constructs a clean representation of the scene background by aggregating static elements across multiple frames. This background reconstruction enables the client to focus on rendering dynamic content without redundant transmission of unchanging or non-salient visual elements.
Once the objects have been extracted, their representations are encoded into a compact format for transmission. The encoding is designed to balance efficiency and fidelity, ensuring that only essential information is relayed while preserving the integrity of motion and appearance. The system employs a structured format that includes object bounding regions, key spatial features, and temporal consistency constraints.

\subsection{PointStream - Client side}
On the client side, the received representations are processed to synthesize the missing visual information and reconstruct a complete scene. A generative model is utilized to infer fine details, reconstruct occluded regions, and blend dynamic objects seamlessly into the scene. This allows the client to render a visually coherent sequence while significantly reducing bandwidth.
The division of labor between server and client ensures that the system can scale across different scene types while maintaining high fidelity and efficiency.

\newpage
\section{Performance Evaluation}

\subsection{PointStream Implementation}
To validate our approach, we applied the system to the analysis of tennis matches. Tennis serves as an ideal test case due to its structured environment, where the primary dynamic elements—players, rackets, and balls—move against a largely static background. This setup allows for efficient segmentation and reconstruction of relevant objects while minimizing interference from extraneous scene elements.
Our implementation is based on a dataset of multiple tennis videos from top-level tournaments. The processing pipeline, implemented in Python, manages video processing, object extraction, background reconstruction, keypoint extraction, and optionally object segmentation (for model training purposes). The entire system is packaged into a Docker container and deployed on a local Ubuntu 20.04 LTS server, equipped with an Intel Xeon Gold 5218 CPU (64 cores, 2.30 GHz) and an NVIDIA Quadro RTX 8000 GPU with 48 GB of memory.

\subsubsection{Video Processing and Object Detection}
The system processes input videos by extracting frames and applying YOLO, an object detection model, to identify relevant objects within each scene. To capture fast-moving elements prone to motion blur, YOLO is configured with a confidence threshold of 0.25 and an IoU threshold of 0.1—values lower than typical settings. While this improves detection recall, it also increases the number of recognized elements, leading to higher computational demands. Despite YOLO’s efficiency, handling an excessive number of detections presents a challenge, requiring an optimized selection mechanism.

\subsubsection{Object Association}
To streamline processing, only objects of interest—players, rackets, and balls—are retained, while all other detections are discarded as background. Since many people appear in the scene but do not participate in gameplay, further refinement is necessary to distinguish between players and bystanders. This is achieved by associating rackets with players based on the spatial overlap of their bounding boxes.
Rackets, due to their rapid motion, are often detected inconsistently across frames. To address this, once a racket is linked to a player, the system continuously tracks that player, dynamically adjusting their bounding box in subsequent frames to ensure their racket remains within the detection region—even if temporarily unrecognized.

\subsubsection{Data Storage}
Once objects have been categorized and associated, they are stored systematically. Each detected instance is cropped and saved into object-specific directories, while bounding box coordinates are recorded for later use and transmission. 
To ensure robustness, objects that appear inconsistently across frames—such as bystanders or erroneous detections—are filtered out.
This structured storage facilitates downstream processing, ensuring efficient access to detection data for subsequent stages of the system.

\subsubsection{Segmentation and Final Processing}
After detecting and filtering objects, the system applies a second, YOLO-based, processing stage focused on pose estimation and, in case of training, segmentation. Since bounding boxes have already been extracted, this step can be performed in parallel for each object, significantly improving speed.
The result of this step will be a number of keypoints for each detected object, capturing essential structural information, such as player poses and racket orientations. 
During training, the system will also generate masks that segment each salient object in the image. For the tennis task at hand, these masks typically include a ball, and/or a player-racket pair. Each class of object is masked with a different color, to allow the client to correctly identify and train on different objects separately, and regenerate them with higher fidelity.

\subsubsection{Background Reconstruction}
Once all frames and objects have been processed, the system reconstructs the background by aggregating multiple frames that had their dynamic objects obscured. In other words, regions of the background that were occluded by the moving objects, are filled with information from other frames where these objects are in a different location.
The extracted segmentation masks and keypoints are stored alongside their corresponding bounding box data. This structured dataset allows the system to learn meaningful object representations, improving the accuracy of generative models used in later stages.

\subsubsection{Data Export}
Finally, the extracted bounding box data is saved to a CSV file, containing frame numbers, object IDs, class labels, and keypoint coordinates where applicable. It can then be broadcast to the client, at a much lower bitrate than a traditionally-encoded video at similar resolution.

\subsubsection{Client-Side Processing and Reconstruction}


\subsection{Experimental Settings}
\textbf{Encoding}.

\textbf{Content}.

\textbf{Training and testing settings.}

\textbf{Scenarios}.

\textbf{Evaluation metrics}.
\subsection{Results and Analysis}
% \subsubsection{Quantitative and Qualitative Comparisons}
We conduct a comprehensive comparison between our PointStream and
several state-of-the-art X methods, including X, Y, Z.

\noindent

\textbf{Quantitative comparison.} Table with results about PSNR etc.

\noindent

\textbf{Qualitative comparison.} Exemplar visual results for various methods are presented in Figure X. PointStream generates high-quality...

% \noindent
% \textbf{Dataset.}

% \noindent

% \noindent
% \textbf{Evaluation metrics.}

% \noindent
% \textbf{Implementation details.}



% \subsection{Ablation Study}

% Ablation study here.

\section{Conclusion}

\newpage

\section{Acknowledgments}
The financial support of the Austrian Federal Ministry for Digital and Economic Affairs, the National Foundation for Research, Technology and Development, and the Christian Doppler Research Association, is gratefully acknowledged. Christian Doppler Laboratory ATHENA: https://athena.itec.aau.at/

% The next two lines define the bibliography style to be used, and the bibliography file.
%\printbibliography

\bibliographystyle{acm}
\bibliography{ref}

\end{document}